\subsection{The actual transfer at an arbitrary partial head}
\label{actual-partial-head-transfer}

The previous form inequalities apply to a partial-head coefficient
only after its integral and its next-level coefficient have been
identified.  We establish that identification below every terminal
level.  A head can contain blocks of any positive sizes; its internal
leading minors are allowed to vanish.

Fix the original normalized tower, its compact probability measures,
its recursively defined unipotent Haar measures, and a compact smooth
coefficient $F$ with positive Gram cutoff.  Write
\[
 h=((s_1,t_1),\ldots,(s_a,t_a)),\qquad
 j=\sum_b s_b\le r,\qquad t_b\ne0,
 \qquad Z_h=\operatorname{diag}(t_1J_{s_1},\ldots,t_aJ_{s_a}),
\]
where $J_s$ reverses the $s$ coordinates.  The source slice
$S_hF$ is the integral \texttt{sliceAt}; it uses the actual
Lebesgue root coordinates, phase and block matrix
\texttt{PhiFun}, not a different orbit measure.
For an original signed word $w$ and a component $e$, put
\[
       F_{j,w,h,e}=S_h(\mathsf V_{j,w,e}F).
\]
Here $\mathsf V_{j,w,e}$ is exactly the tail-first composition
\texttt{uwG} of the normalized isotropic Capelli letters.
At $j=0$ its component with a letter index $1$ is zero,
as in the original definition.

\begin{theorem}\label{partial-head-transfer-actual}
Every $F_{j,w,h,e}$ is a compact smooth coefficient on $G_{r-j}$
with a positive Gram cutoff.  These coefficients and all their
parameter derivatives have common compact support and a common
positive cutoff when the nonzero head parameters range in a compact
set of a fixed block pattern.

If $j<r$, adjoining a singleton $(1,c)$, $c\ne0$, transfers the
Jacquet transforms of these vectors by the original one-coordinate
word operator.  In particular, for all heads $h,h'$ of size $j$
and every original admissible finite signed word $w$, the actual
radial marginal at $c$ equals
\begin{equation}\label{partial-head-density}
 |c|^{\,2(2(r-j)+k-1)}
    \operatorname{RR}_{j+1,w}
          (h\mathbin{\|}(1,c),h'\mathbin{\|}(1,c)).
\end{equation}
It has the continuous integrable extension and Fourier identity
provided by the established one-coordinate radial theorem.
This supplies the original compact-slice, head-Jacquet and
below-terminal radial-head interfaces, with no new norm constant.
\end{theorem}

\paragraph{The quotient coordinates and their properness.}
Let $p$ be the size of a fixed head and let $\rho=r-p$.
In the ordered decomposition of dimensions $(p,2\rho+k+1,p)$,
write $J=J_p$ and let $H$ be the middle Hermitian matrix.  The
original radical coordinate is
\[
 R(X,K)=
 \begin{pmatrix}
 I&-JX^*H&J(-\tfrac12X^*HX+K)\\
 0&I&X\\0&0&I
 \end{pmatrix},\qquad K^*=-K.
\]
Its coordinates are the complex entries of $X$, the strict upper
complex entries of $K$, and the real numbers $-iK_{aa}$.
Their product Lebesgue measure is the prescribed one.
For $\det Z\ne0$ set
\[
 m(Z,\gamma)=
 \begin{pmatrix}0&0&(Z^*)^{-1}J\\0&\gamma&0\\JZ&0&0\end{pmatrix}.
\]
For $g=R(X_L,K_L)^{-1}m(Z,\gamma)R(X_R,K_R)$, put
$D=g_{31}=JZ$.  Block multiplication gives
\begin{equation}\label{partial-head-inverse-coordinates}
 \begin{aligned}
 X_L&=-g_{21}D^{-1},&
 X_R^*&=-JD^{-1}g_{32}H^{-1},\\
 K_L&=-Jg_{11}D^{-1}-\tfrac12X_L^*HX_L,&
 K_R&=JD^{-1}g_{33}+\tfrac12X_R^*HX_R,\\
 \gamma&=g_{22}-g_{21}D^{-1}g_{32}.&&
 \end{aligned}
\end{equation}
Thus these coordinates are recovered rationally from $g$, with
denominators only powers of $\det Z$ and its conjugate.

For the head itself write $Z=LZ_hU$, where $L$ is block lower
unitriangular for $(s_1,\ldots,s_a)$, and $U$ is upper
unitriangular.  This is precisely the source's restricted
lower-root chart: conjugation by the block reversing permutation
makes an internal lower root upper, so it is omitted from $L$.
Block Gaussian elimination recovers $L$ and $B=Z_hU$ from $Z$:
\[
 B_{bc}=Z_{bc}-\sum_{d<b}L_{bd}B_{dc}\quad(c\ge b),\qquad
 L_{cb}=\left(Z_{cb}-\sum_{d<b}L_{cd}B_{db}\right)B_{bb}^{-1}
       \quad(c>b).
\]
Here $\det B_{bb}=\det(t_bJ_{s_b})\ne0$.  Adjugates therefore
give polynomial bounds for $L,B,U$ in the entries of $Z$, locally
uniformly in the nonzero $t_b$.  Internal principal minors inside
a block are never divided by.  Together with
\eqref{partial-head-inverse-coordinates}, this proves that the
preimage of a compact set of group matrices is compact in all
the quotient coordinates and in $\gamma$.

The omitted roots are essential.  In a block of size two,
\[
 \begin{pmatrix}1&0\\a&1\end{pmatrix}
 J_2
 \begin{pmatrix}1&-a\\0&1\end{pmatrix}=J_2
 \qquad(a\in\mathbb C).
\]
Integrating both full triangular groups would introduce an
unbounded stabilizer fibre, even over a compact coefficient.
The restricted chart removes exactly this fibre.  Its use does
not replace or renormalize the original stabilizer quotient.

The same calculation with a final unfixed invertible block $Z_0$
uses the Schur coordinates
\[
 Z=\begin{pmatrix}I&0\\ A&I\end{pmatrix}
       \begin{pmatrix}H_0&0\\0&Z_0\end{pmatrix}
       \begin{pmatrix}I&B\\0&I\end{pmatrix},
 \quad A=Z_{21}H_0^{-1},\quad B=H_0^{-1}Z_{12}.
\]
On a compact subset of $\det Z_0\ne0$ all inverse bounds remain
valid.  The source phase has modulus one.  Hence differentiation
under the partial-head integral is justified by a common compact
support in its integration variables.  The Gram identity is
\begin{equation}\label{partial-head-gram}
 \mathcal G_r\bigl(R_L^{-1}m(Z,\gamma)R_R\bigr)
                        =|\det Z|^2\mathcal G_\rho(\gamma).
\end{equation}
Indeed each larger flag minor is $\det Z$ times its smaller
counterpart; the two terms defining the full Gram polynomial
have the same factor.  Thus the cutoff is
$\varepsilon/|\det Z_h|^2$, with its minimum over the parameter
compact set when a family is considered.
The operators $\mathsf V_{j,w,e}$ have coefficients smooth on
$D_j\ne0$ and do not enlarge support.  On the integration chart
$D_j=\det Z_h$ is fixed and nonzero.  The same compact-support
argument applies to every derivative of every finite word.
This proves the first assertion of the theorem, including $p=0$.

\paragraph{Adjoining a singleton, with the original measure.}
The coordinate factorization can be done before taking any
quadratic form.  Denote the head permutation by $\sigma$ and write
its original matrix as $H_0=w_\sigma L_0D_0U_0$.
The new lower and upper off-diagonal blocks, initially $V$ and $B_0$,
are changed to
\[
                  A=VL_0^{-1}w_\sigma^{-1},\qquad
                  B=U_0^{-1}B_0.
\]
Their real Jacobian is one: $L_0,U_0$ are unitriangular and
$w_\sigma$ is a permutation.  The resulting head matrix is
\[
 \begin{pmatrix}I&0\\ A&1\end{pmatrix}
       \operatorname{diag}(H_0,c)
       \begin{pmatrix}I&B\\0&1\end{pmatrix}.
\]
The lower subdiagonal and upper superdiagonal sums in the phase
split into the old sums plus $A_{p-1}$ and $B_{p-1}$.
For example the last column of the old lower unitriangular
matrix is the last coordinate vector, so the new subdiagonal
entry of $A L$ is exactly $A_{p-1}$.  This verifies the phase,
including the two corner terms, not just its modulus.

Splitting the radical first at $p$ and then at one paired vector
gives the same coordinates as splitting first at $p+1$ and then
making the displayed shear change.  This equality follows by
multiplying the three displayed block matrices $R,m,R$.
Its Haar Jacobian is also one.  One way to see the normalization
directly is to order the original root coordinates by height:
the changed coordinate for each root is its old coordinate plus
a polynomial in lower-height coordinates.  The possible signs
and coordinate permutations come from $J$ and $w_\sigma$ and
have real Jacobian of absolute value one.  The central terms
$-X^*HX/2$ are translations in the central coordinates.  The
restriction to each new radical is consequently exactly its
Lebesgue Haar measure in the recursion defining $\mu_N$.
Compact conjugation has determinant of absolute value one and
preserves the restricted characters; its probability measures
agree under the tower embeddings.  Fubini is legitimate by the
proved properness.  We obtain, for the empty word,
\begin{equation}\label{partial-head-empty-step}
 J_{r-j-1}(S_{h\mathbin{\|}(1,c)}F)(\gamma)
       =J_{r-j}(S_hF)(m(c,\gamma)),\qquad c\ne0.
\end{equation}
No power of $c$ occurs in this ungauged identity.

\paragraph{The Capelli calculation at an invertible head.}
The word identity requires more than
\eqref{partial-head-empty-step}.  Put
$\kappa=(r-j)+k/2$.  At
$g=m(\operatorname{diag}(Z,c),\gamma)$, where
$Z\in\operatorname{GL}_j(\mathbb C)$, the relevant equivariant
coefficient has the right upper-root character law and the left
character law for the roots joining the first $j$ vectors to
the next one.  If $\Lambda(X)=(XF)(g)$, these laws give
\[
 \Lambda(XE_{i,j+1})=\delta_{ij}\Lambda(X),\qquad
 \Lambda(E_{j+1,b}X)=c(Z^{-1})_{bj}\Lambda(X).
\]
Indices here start at one.  The second equality holds after
every word of right derivatives, since those derivatives commute
with left translations.  The two real torus directions give
$E_{j+1,j+1}=c\partial_c$ on this one-parameter restriction.
In the dual basis their values are the conjugate lower-root
scalars and $-\bar c\partial_{\bar c}$.

Apply the established last-border Capelli identity to the first
$j$ indices and the new one.  Only the last right-root scalar
survives.  The sum of the remaining cofactors is
\[
 \sum_{b=1}^j(-1)^{j+b}(Z^{-1})_{bj}
       Q_{[j]\setminus j,[j]\setminus b}(u)
       =(\det Z)^{-1}\mathcal A_{j-1}(u),
\]
because
$(-1)^{j+b}\det Z\,(Z^{-1})_{bj}$ is the determinant of
the first $j-1$ rows with column $b$ omitted.  This is the
ordinary adjugate formula; it does not assume that $Z$ is
diagonal.  The full size-$j$ contraction is
$\det Z\,Q_{[j],[j]}(u)$.  Dividing by the actual size-$j$
and size-$(j+1)$ minors gives the plus recurrence.  Apply the
reverse-row border identity in the dual ordering to obtain the
minus recurrence.  With the original spectral shifts they are
\begin{equation}\label{partial-head-letter-recurrence}
 \binom{v^+_{j+1,0}}{v^+_{j+1,1}}
  =\begin{pmatrix}Y+\kappa+c\partial_c&-c\\c^{-1}&0\end{pmatrix}
          \binom{v^+_{j,0}}{v^+_{j,1}},\qquad
 \binom{v^-_{j+1,0}}{v^-_{j+1,1}}
  =\begin{pmatrix}Y-\kappa-\bar c\partial_{\bar c}&-\bar c\\
                   \bar c^{-1}&0\end{pmatrix}
          \binom{v^-_{j,0}}{v^-_{j,1}}.
\end{equation}
The case $j=0$ follows directly from the empty determinant and
the prescribed zero second component.  In the first calculation
$u=Y+(2r+k)/2$, so $u-j=Y+\kappa$; in the second it is
$Y-(2r+k)/2+j=Y-\kappa$.

In this argument the coefficient minors must be frozen at the
point before forming the enveloping-algebra expression.  Their
first $j$ rows at $m(\operatorname{diag}(Z,c),\gamma)$ depend
only on $Z$, so they are constant as $c$ varies.  They and their
dual versions also commute with the new torus generator.
These observations justify replacing that generator by the
parameter derivative in \eqref{partial-head-letter-recurrence}.
They avoid differentiating a coefficient which was meant to
be frozen.

There is one further point needed for word induction.  A size-$j$
letter restricted to this chart differentiates only the head
matrix $Z$.  After evaluating its coefficient minors, all surviving
quantum-minor indices lie in the first $j$ coordinates.  Their
complexified right directions change $Z$ and its Hermitian inverse
block, and leave $c$ and $\gamma$ fixed.  The resulting
$Z$-differential expression is independent of $c$.  It therefore
commutes with multiplication by $c,\bar c$ and with
$\partial_c,\partial_{\bar c}$.  This is the specific commutation
used in the induction; no commutation of two arbitrary head
letters is asserted.

The normalized letters preserve the character laws just used.
Indeed the first $j$ flag vectors are fixed by the relevant
left subgroup, and right upper-unipotent translation changes
the contracted rows and quantum columns by inverse triangular
operations.  Exterior covariance cancels these operations,
while $D_j$ is invariant.  The same argument applies to the
dual contraction.  Thus it applies after every earlier letter.
All partial averages in the preceding paragraph commute with
these covariant operators.  The common compact majorants justify
moving each derivative through the integral.  Induction in the
original tail-first order therefore proves the word form of
\eqref{partial-head-letter-recurrence}, with a fresh auxiliary
index at each step.  This is exactly the original one-coordinate
head-Jacquet operator.  No two noncommuting letter operators
are interchanged.

For clarity, the source's \texttt{oneBT} is the matrix action of
\texttt{bT} at the site with operators
$c\partial_c+\kappa$, $-(\bar c\partial_{\bar c}+\kappa)$,
$c,\bar c,c^{-1},\bar c^{-1}$.
Its defining recursion is $bT(\varnothing)=1$ and
$bT(l::w)=T_l\,bT(w)$, with $T_l$ inserted at the tensor position
of $l$.  This is the same head-after-tail order as
\texttt{uwG}; no rational $R$-matrix dressing enters this operator.
If two punctured-smooth representatives agree on $c\ne0$, all
their local derivatives there agree.  Every finite product just
displayed therefore gives the same values there, irrespective
of the representatives' values at zero.  This verifies the
representative quantifier in the original head-Jacquet interface.

\paragraph{The density and Fourier transfer.}
The radial gauge at the step is
$\lambda(c)=|c|^{2(r-j)+k}=|c|^{2\kappa}$.
Conjugating \eqref{partial-head-letter-recurrence} by $\lambda$
removes $\kappa$ and gives precisely
\[
       \begin{pmatrix}Y+c\partial_c&-c\\c^{-1}&0\end{pmatrix},
       \qquad
       \begin{pmatrix}Y-\bar c\partial_{\bar c}&-\bar c\\
                         \bar c^{-1}&0\end{pmatrix}.
\]
The existing residue-congruence and integrated-descent calculation
now applies to the actual compact vectors $F_{j,w,h,e}$.
Its measure in this coordinate is $|c|^{-2}dc$, and hence its
ungauged scalar density is exactly
\[
             |c|^{-2}\lambda(c)^2
                      =|c|^{2(2(r-j)+k-1)}.
\]
The normalized terminal Haar measure and the tower give the same
positive Jacquet factor on both sides; equivalently the recursive
one-pair integration formula supplies this equality before a
constant is introduced.  Thus no unknown ratio of level constants
is used in \eqref{partial-head-density}.

The established Schwartz-marginal theorem, including its finite-word
version, supplies a Schwartz extension of this density and the
Fourier pairing with the actual wave $\exp(sc-\bar s\bar c)$.
At every $c\ne0$ the just-proved compact-family regularity identifies
the density with that extension, not merely almost everywhere.
For two different heads the same argument is sesquilinear and
uses a common minimum of their positive cutoffs.
This proves the theorem \cite{partial-head-source,partial-head-established}.

The terminal comparison is deliberately separate.  For $j=r$ the
next expression is built from the original top lift and the root
quotient orbital vector, rather than a lower-rank residue form.
The top-lift hypothesis identifies its action on the regular radial
cell.  Extending that identity to an arbitrary block head, with
the prescribed stabilizer measure, is the remaining terminal
compatibility statement.  The theorem above does not assume it
and does not declare the complete all-level kernel estimate proved.
